\begin{xltabular}{\textwidth}{|>{\raggedright\arraybackslash}p{3cm}|>{\raggedright\arraybackslash}p{4cm}|X|}
    \caption{Kriteria Evaluasi Infrastruktur Penyimpanan Audit Trail}
    \label{tab:storage-criteria} \\
    \hline
    \textbf{Kriteria} &
    \textbf{Deskripsi Singkat} &
    \textbf{Justifikasi Pemilihan} \\
    \hline
    \endfirsthead

    \hline
    \textbf{Kriteria} &
    \textbf{Deskripsi Singkat} &
    \textbf{Justifikasi Pemilihan} \\
    \hline
    \endhead

    Immutability &
    Kemampuan teknologi penyimpanan untuk mencegah data yang telah disimpan dimodifikasi atau dihapus. &
    Audit trail berfungsi sebagai bukti aktivitas sistem sehingga data yang telah disimpan harus terlindungi dari perubahan maupun penghapusan. \\
    \hline

    Durability &
    Kemampuan teknologi penyimpanan untuk mempertahankan data dari risiko kehilangan akibat kegagalan perangkat keras, kerusakan media, atau gangguan sistem. &
    Audit trail umumnya disimpan dalam jangka waktu yang panjang sehingga kehilangan data dapat menghilangkan bukti penting pada proses audit maupun investigasi. \\
    \hline

    Availability &
    Kemampuan teknologi penyimpanan untuk menyediakan akses terhadap data ketika dibutuhkan oleh pengguna atau sistem yang berwenang. &
    Audit trail harus dapat diakses sewaktu-waktu, terutama ketika dilakukan proses audit, investigasi insiden, atau analisis forensik. \\
    \hline

    Scalability &
    Kemampuan teknologi penyimpanan untuk menangani pertumbuhan volume data tanpa memerlukan perubahan arsitektur yang signifikan. &
    Audit trail dihasilkan secara terus-menerus selama sistem beroperasi sehingga volumenya akan terus bertambah. \\
    \hline

    Retention Capability &
    Kemampuan teknologi penyimpanan untuk menerapkan kebijakan retensi data sesuai periode penyimpanan yang ditentukan. &
    Banyak organisasi maupun regulasi mewajibkan audit trail disimpan dalam jangka waktu tertentu. Dukungan terhadap kebijakan retensi mempermudah pengelolaan siklus hidup audit trail sekaligus membantu memenuhi kebutuhan tata kelola dan kepatuhan. \\
    \hline

    Ease of Integration &
    Kemudahan teknologi penyimpanan untuk diintegrasikan dengan komponen lain dalam arsitektur sistem, seperti aplikasi, layanan API, maupun mekanisme keamanan tambahan. &
    Penelitian ini berfokus pada perancangan solusi sehingga teknologi penyimpanan yang dipilih perlu dapat diintegrasikan dengan arsitektur yang diusulkan tanpa memerlukan kompleksitas implementasi yang berlebihan. Kemudahan integrasi juga memengaruhi efisiensi pengembangan dan pemeliharaan sistem. \\
    \hline

\end{xltabular}